
\documentclass{article}
\usepackage{colortbl}
\usepackage{makecell}
\usepackage{multirow}
\usepackage{supertabular}

\begin{document}

\newcounter{utterance}

\centering \large Interaction Transcript for game `matchit\_ascii', experiment `similar\_grid\_1', episode 0 with NEW\_Reflect\_Agent.
\vspace{24pt}

{ \footnotesize  \setcounter{utterance}{1}
\setlength{\tabcolsep}{0pt}
\begin{supertabular}{c@{$\;$}|p{.15\linewidth}@{}p{.15\linewidth}p{.15\linewidth}p{.15\linewidth}p{.15\linewidth}p{.15\linewidth}}
    \# & $\;$A & \multicolumn{4}{c}{Game Master} & $\;\:$B\\
    \hline

    \theutterance \stepcounter{utterance}  
    & & \multicolumn{4}{p{0.6\linewidth}}{
        \cellcolor[rgb]{0.9,0.9,0.9}{
            \makecell[{{p{\linewidth}}}]{
                \texttt{\tiny{[P1$\langle$GM]}}
                \texttt{You are participating in a collaborative guessing game.} \\
\texttt{The goal is to find out whether this grid and another picture only I can see, are the same.} \\
\texttt{Please describe your grid first. Then, I will provide my description and we can ask each other questions about the grids to figure out whether they are the same. Now start your short grid description with "DESCRIPTION:" followed by the description. Do not add anything else. This is your grid:} \\
\texttt{X ▢ ▢ ▢ ▢} \\
\texttt{X X ▢ ▢ ▢} \\
\texttt{X X X ▢ ▢} \\
\texttt{X X X X ▢} \\
\texttt{X X X X X} \\
            }
        }
    }
    & \\ \\

    \theutterance \stepcounter{utterance}  
    & \multicolumn{4}{p{0.45\linewidth}}{
        \cellcolor[rgb]{0.95,0.95,0.95}{
            \makecell[{{p{\linewidth}}}]{
                \texttt{\tiny{[P1$|$internal]}}
                \texttt{Reflections from previous episodes:} \\
\texttt{{-}} \\
\texttt{ Failed. I described the grid as a symmetrical X pattern, which was correct, but I didn’t ask clarifying questions about the other player’s grid to confirm details like exact dimensions or symmetry. Next time, I would follow my description with a direct question about their grid’s structure to compare systematically.} \\
\texttt{{-}} \\
\texttt{ Failed. I gave a vague description ("grid with X's forming patterns, mostly in center columns") that didn't capture the exact shape. My partner's description didn't match, but I couldn't pinpoint differences because my initial description lacked precise row{-}by{-}row or coordinate details. Next time, I would give a structured description (e.g., row{-}by{-}row counts or positions) to enable exact comparison.} \\
\texttt{{-}} \\
\texttt{ Failed. I described the grid as a cross shape, but the user's hidden grid was a plus sign. The difference was orientation: mine had a vertical line of five with a horizontal line intersecting at the center, forming a symmetrical cross. The user's plus likely had a different central alignment. I should have asked a clarifying question about the exact position and length of the lines instead of assuming "cross" was unambiguous. Next time, specify the row/column coordinates of the lines.} \\
            }
        }
    }
    & & \\ \\

    \theutterance \stepcounter{utterance}  
    & \multicolumn{4}{p{0.45\linewidth}}{
        \cellcolor[rgb]{0.95,0.95,0.95}{
            \makecell[{{p{\linewidth}}}]{
                \texttt{\tiny{[P1$|$internal]}}
                \texttt{{[}attempt 1 – valid{]}} \\
\texttt{DESCRIPTION: A 5x5 grid where the X's form a diagonal pattern starting from the top{-}left corner, with each row having one more X than the previous row, filling the leftmost columns. Specifically: Row 1 has an X in column 1, Row 2 has X's in columns 1{-}2, Row 3 has X's in columns 1{-}3, Row 4 has X's in columns 1{-}4, and Row 5 has X's in all five columns. The remaining cells are empty (▢).} \\
            }
        }
    }
    & & \\ \\

    \theutterance \stepcounter{utterance}  
    & \multicolumn{4}{p{0.6\linewidth}}{
        \cellcolor[rgb]{0.9,0.9,0.9}{
            \makecell[{{p{\linewidth}}}]{
                \texttt{\tiny{[P1$\rangle$GM]}}
                \texttt{DESCRIPTION: A 5x5 grid where the X's form a diagonal pattern starting from the top{-}left corner, with each row having one more X than the previous row, filling the leftmost columns. Specifically: Row 1 has an X in column 1, Row 2 has X's in columns 1{-}2, Row 3 has X's in columns 1{-}3, Row 4 has X's in columns 1{-}4, and Row 5 has X's in all five columns. The remaining cells are empty (▢).} \\
            }
        }
    }
    & & \\ \\

    \theutterance \stepcounter{utterance}  
    & & & \multicolumn{2}{p{0.3\linewidth}}{
        \cellcolor[rgb]{0.9,0.9,0.9}{
            \makecell[{{p{\linewidth}}}]{
                \texttt{\tiny{[GM$|$GM]}}
                \texttt{continue} \\
            }
        }
    }
    & & \\ \\

    \theutterance \stepcounter{utterance}  
    & & \multicolumn{4}{p{0.6\linewidth}}{
        \cellcolor[rgb]{0.9,0.9,0.9}{
            \makecell[{{p{\linewidth}}}]{
                \texttt{\tiny{[GM$\rangle$P2]}}
                \texttt{You are participating in a collaborative guessing game.} \\
\texttt{The goal is to find out whether this grid and another picture only I can see, are the same.} \\
\texttt{Please describe your grid first. Then, I will provide my description and we can ask each other questions about the grids to figure out whether they are the same. Now start your short grid description with "DESCRIPTION:" followed by the description. Do not add anything else. This is your grid:} \\
\texttt{▢ ▢ ▢ ▢ X} \\
\texttt{▢ ▢ ▢ X X} \\
\texttt{▢ ▢ X X X} \\
\texttt{▢ X X X X} \\
\texttt{X X X X X} \\
            }
        }
    }
    & \\ \\

    \theutterance \stepcounter{utterance}  
    & & & \multicolumn{4}{p{0.45\linewidth}}{
        \cellcolor[rgb]{0.95,0.95,0.95}{
            \makecell[{{p{\linewidth}}}]{
                \texttt{\tiny{[P2$|$internal]}}
                \texttt{Reflections from previous episodes:} \\
\texttt{{-}} \\
\texttt{ Failed. My description was too vague ("symmetric pattern," "T shape," "U shape") and could match multiple grids. I should have given a precise structural description, like "X's are on all positions where row index equals column index or row index + column index equals 4" (the main diagonals). Next time, I will use exact mathematical or coordinate{-}based descriptions to avoid ambiguity.} \\
\texttt{{-}} \\
\texttt{ Failed. My description was inaccurate: I described a diagonal pattern, but my grid actually had a vertical "X" band in columns 2{-}4, not a true diagonal. This misled the comparison from the start. Next time, I would describe the pattern more precisely: "A 5x5 grid with a vertical block of X's in columns 2, 3, and 4, except for some missing X's forming a checkerboard{-}like absence in the center rows."} \\
\texttt{{-}} \\
\texttt{ FAILED. My initial description was too brief, lacking precise coordinates, which likely caused confusion. The partner's detailed correction suggests mismatched interpretation. Next time, I will provide a full coordinate{-}based description immediately to avoid ambiguity and align understanding from the start.} \\
            }
        }
    }
     \\ \\

    \theutterance \stepcounter{utterance}  
    & & & \multicolumn{4}{p{0.45\linewidth}}{
        \cellcolor[rgb]{0.95,0.95,0.95}{
            \makecell[{{p{\linewidth}}}]{
                \texttt{\tiny{[P2$|$internal]}}
                \texttt{{[}attempt 1 – valid{]}} \\
\texttt{DESCRIPTION: A 5x5 grid where X's are present in all positions (i, j) such that i + j ≥ 4, with i and j ranging from 0 to 4 (0{-}indexed).} \\
            }
        }
    }
     \\ \\

    \theutterance \stepcounter{utterance}  
    & & & \multicolumn{4}{p{0.6\linewidth}}{
        \cellcolor[rgb]{0.9,0.9,0.9}{
            \makecell[{{p{\linewidth}}}]{
                \texttt{\tiny{[GM$\langle$P2]}}
                \texttt{DESCRIPTION: A 5x5 grid where X's are present in all positions (i, j) such that i + j ≥ 4, with i and j ranging from 0 to 4 (0{-}indexed).} \\
            }
        }
    }
     \\ \\

    \theutterance \stepcounter{utterance}  
    & & & \multicolumn{2}{p{0.3\linewidth}}{
        \cellcolor[rgb]{0.9,0.9,0.9}{
            \makecell[{{p{\linewidth}}}]{
                \texttt{\tiny{[GM$|$GM]}}
                \texttt{continue} \\
            }
        }
    }
    & & \\ \\

    \theutterance \stepcounter{utterance}  
    & & \multicolumn{4}{p{0.6\linewidth}}{
        \cellcolor[rgb]{0.9,0.9,0.9}{
            \makecell[{{p{\linewidth}}}]{
                \texttt{\tiny{[GM$\rangle$P2]}}
                \texttt{This is my DESCRIPTION: A 5x5 grid where the X's form a diagonal pattern starting from the top{-}left corner, with each row having one more X than the previous row, filling the leftmost columns. Specifically: Row 1 has an X in column 1, Row 2 has X's in columns 1{-}2, Row 3 has X's in columns 1{-}3, Row 4 has X's in columns 1{-}4, and Row 5 has X's in all five columns. The remaining cells are empty (▢).} \\
\texttt{Now ask a question in order to find out new aspects of my grid that may be different to your grid. Start with "QUESTION:" and do not add anything else.} \\
            }
        }
    }
    & \\ \\

    \theutterance \stepcounter{utterance}  
    & & & \multicolumn{4}{p{0.6\linewidth}}{
        \cellcolor[rgb]{0.9,0.9,0.9}{
            \makecell[{{p{\linewidth}}}]{
                \texttt{\tiny{[GM$\langle$P2]}}
                \texttt{DESCRIPTION: A 5x5 grid where the X's form a diagonal pattern starting from the top{-}right corner, with each row having one more X than the previous row, filling the rightmost columns. Specifically: Row 1 has an X in column 5, Row 2 has X's in columns 4{-}5, Row 3 has X's in columns 3{-}5, Row 4 has X's in columns 2{-}5, and Row 5 has X's in all five columns. The remaining cells are empty (▢).  } \\
\texttt{QUESTION: Does your grid have an X in the top{-}right corner (row 1, column 5)?} \\
            }
        }
    }
     \\ \\

    \theutterance \stepcounter{utterance}  
    & & & \multicolumn{2}{p{0.3\linewidth}}{
        \cellcolor[rgb]{0.9,0.9,0.9}{
            \makecell[{{p{\linewidth}}}]{
                \texttt{\tiny{[GM$|$GM]}}
                \texttt{abort, first word: DESCRIPTION:, but expected QUESTION:} \\
            }
        }
    }
    & & \\ \\

\end{supertabular}
}

\end{document}
